% Gummel plot (bjt-gummel): the set-up a person would wire, semantically the deck's.
%   BG 0 bg I=10^V(x)   the drive, a log sweep into the base node               -> current source I_source (set)
%   VG cg bg 0          a 0 V source from collector to base: VCB = 0, reads IC  -> ammeter I_C (measured)
%   QG cg bg 0 DUT      the transistor, emitter grounded                        -> npn
%   V(bg) read: VBE                                                             -> voltmeter V_BE (measured)
%   The dial V(x) is the deck's way of sweeping; the drawing names the current it sets.
\documentclass[tikz, border=10pt]{standalone}
\usepackage{bench}
\begin{document}
\begin{tikzpicture}
\benchtitle{0,7.0}{Gummel plot}
\def\gnd{1.0}
% the transistor
\node[npn] (Q) at (4,3) {};
\node[anchor=west, font=\small, align=left] at ($(Q.C)+(0.15,-0.55)$) {DUT};
% emitter to ground
\draw (Q.E) -- (Q.E |- 0,\gnd) node[ground]{};
\nodename{anchor=west}{$(Q.E)+(0.1,-0.3)$}{0}
% the ammeter in the collector, up to the rail
\coordinate (ctop) at ($(Q.C)+(0,1.7)$);
\ammeter{Q.C}{ctop}{$I_C$}
\nodename{anchor=west}{$(Q.C)+(0.1,0.25)$}{cg}
% the base: a current source from ground
\coordinate (b) at (Q.B);
\coordinate (bsrc) at ($(b)+(-1.4,0)$);
\draw (b) -- (bsrc);
\draw[set] (bsrc |- 0,\gnd) node[ground, black]{} to[I, l_=$I_{source}$] (bsrc);
\nodename{anchor=south}{$(b)+(-0.5,0.02)$}{bg}
% the voltmeter from the base wire to ground
\coordinate (bvm) at ($(b)+(-3.2,0)$);
\draw (b) -- (bvm);
\voltmeter{bvm |- 0,\gnd}{bvm}{$V_{BE}$}
\draw (bvm |- 0,\gnd) node[ground]{};
% the loop from the ammeter back to the base node: VCB = 0
\coordinate (bj) at ($(b)+(-0.9,0)$);
\draw (ctop) -| (bj);
% the definitions
\definitions{0,-0.1}{%
  $I_{source} = 10^{k}\,\mathrm{A}$, $k = -9,\,-8.95,\,\dots,\,0$ \quad (20 points per decade, 1 nA to 1 A)\\[2pt]
  $V_{CB} = 0$: the ammeter drops nothing, so $V_{CE} = V_{BE}$\\[2pt]
  read $I_C$ and $V_{BE}$;\quad $I_B = I_{source} - I_C$;\quad $h_{FE} = I_C / I_B$%ENDDEFS
}
\end{tikzpicture}
\end{document}
